\chapter{Constructions and enveloping algebras}
\label{ch:constructions}

\section{Brackets, quotients, and representations}

We work over a field \(k\). Statements involving factorials or exponential
series will explicitly assume characteristic zero. All tensor products
in this chapter are over \(k\). An associative algebra is unital, and a
module over it is a unital left module.

\begin{definition}[Lie algebras and their quotients]\label{res:A01}
A Lie algebra is a vector space \(\mathfrak g\) with a bilinear operation
\([x,y]\) satisfying
\[
 [x,x]=0,\qquad
 [x,[y,z]]+[y,[z,x]]+[z,[x,y]]=0.
\]
A linear subspace is a subalgebra if it is closed under brackets; it is
an ideal \(I\) if \([\mathfrak g,I]\subseteq I\). A homomorphism preserves
brackets. For an ideal, the quotient has the bracket
\([x+I,y+I]=[x,y]+I\). If \(f:\mathfrak g\to\mathfrak h\) is a
homomorphism, it induces an isomorphism
\(\mathfrak g/\ker f\longrightarrow\operatorname{im}f\).
\end{definition}

\begin{proof}
Bilinearity and \([x+y,x+y]=0\) imply \([x,y]=-[y,x]\), including in
characteristic two; the alternating identity itself remains part of the
definition. If \(x\) and \(y\) are replaced by \(x+a,y+b\), with
\(a,b\in I\), the bracket changes by
\([a,y]+[x,b]+[a,b]\in I\). Thus the quotient bracket is well defined.
Its alternating and Jacobi identities are images of the identities
in \(\mathfrak g\). The kernel of a homomorphism is an ideal, its image
is a subalgebra, and \(x+\ker f\mapsto f(x)\) is well defined, injective,
surjective onto the image, and bracket preserving.
\end{proof}

An associative algebra \(A\) becomes a Lie algebra under
\([a,b]=ab-ba\). Expansion gives its Jacobi identity: each of the six
ordered triple products occurs once with each sign. In particular,
\(\mathfrak{gl}(V)=\operatorname{End}_k(V)\) is a Lie algebra.

\begin{proposition}[Central and derived series]\label{res:A02}
Set
\[
 D^0\mathfrak g=\mathfrak g,\quad
 D^{r+1}\mathfrak g=[D^r\mathfrak g,D^r\mathfrak g],\qquad
 \gamma _1\mathfrak g=\mathfrak g,\quad
 \gamma _{r+1}\mathfrak g=[\mathfrak g,\gamma _r\mathfrak g].
\]
Here brackets of subspaces mean their linear span. All these subspaces
are ideals. They define solvability and nilpotence by eventual
vanishing, respectively. Subalgebras and quotients of a solvable
(respectively nilpotent) algebra have the same property. An extension
of solvable algebras is solvable.

The upper central series \(Z_0=0\),
\(Z_{r+1}/Z_r=Z(\mathfrak g/Z_r)\) reaches \(\mathfrak g\) if and only
if \(\mathfrak g\) is nilpotent. Moreover,
\[
 [\gamma_i\mathfrak g,\gamma_j\mathfrak g]
 \subseteq\gamma_{i+j}\mathfrak g.
\]
\end{proposition}

\begin{proof}
Jacobi in the form
\([x,[a,b]]=[[x,a],b]+[a,[x,b]]\) shows that the bracket of two
ideals is an ideal; it proves the assertion about both series by
induction. For a homomorphism \(f\), induction gives
\(f(D^r\mathfrak g)=D^r(f\mathfrak g)\) and similarly for \(\gamma_r\).
The analogous inclusions for a subalgebra suffice. If an ideal \(I\)
and the quotient are solvable, some \(D^r\mathfrak g\) lies in \(I\),
and then \(D^{r+s}\mathfrak g\subseteq D^s I\).

For the last inclusion, induct on \(i\). The case \(i=1\) is the
definition. Jacobi gives
\[
 [[\mathfrak g,\gamma_i],\gamma_j]
 \subseteq[\mathfrak g,[\gamma_i,\gamma_j]]
       +[\gamma_i,[\mathfrak g,\gamma_j]]
 \subseteq\gamma_{i+j+1}.
\]
If \(\gamma_{c+1}=0\), induction shows
\(\gamma_{c+1-r}\subseteq Z_r\), starting with
\(\gamma_c\subseteq Z_1\). Thus \(Z_c=\mathfrak g\).
Conversely, \(Z_c=\mathfrak g\) implies
\(\gamma_{r+1}\subseteq Z_{c-r}\), by induction on \(r\), so
\(\gamma_{c+1}=0\).
\end{proof}

Nilpotence implies solvability, since \(D^r\mathfrak g\subseteq
\gamma_{r+1}\mathfrak g\). It does not follow merely from an extension
of two nilpotent algebras: \([h,e]=e\) gives a two-dimensional
counterexample.

\begin{proposition}[Actions, derivations, and semidirect products]
\label{res:A03}
A representation is a linear map
\(\rho:\mathfrak g\to\mathfrak{gl}(M)\) satisfying
\[
 \rho([x,y])=\rho(x)\rho(y)-\rho(y)\rho(x).
\]
The map \(\operatorname{ad}x:y\mapsto[x,y]\) is a representation and
has kernel \(Z(\mathfrak g)\). A derivation of a Lie algebra
\(\mathfrak n\) is a linear map \(d\) with
\(d[a,b]=[da,b]+[a,db]\). Derivations form a Lie subalgebra of
\(\mathfrak{gl}(\mathfrak n)\).

If \(\theta:\mathfrak g\to\operatorname{Der}(\mathfrak n)\) is a
representation, then
\[
 [(x,a),(y,b)]
   =([x,y],\,\theta(x)b-\theta(y)a+[a,b])
 \tag{1.1}\label{eq:semidirect}
\]
makes \(\mathfrak g\oplus\mathfrak n\) a Lie algebra, denoted
\(\mathfrak g\ltimes\mathfrak n\). For an abelian \(\mathfrak n=M\),
the Jacobi identity in this formula is equivalent to the
representation identity.
\end{proposition}

\begin{proof}
Jacobi is exactly the assertion that \(\operatorname{ad}x\) is a
derivation and that \([\operatorname{ad}x,\operatorname{ad}y]
=\operatorname{ad}[x,y]\). Its kernel is the stated center.
Expanding \(de[a,b]-ed[a,b]\) cancels the two mixed terms and proves
that \([d,e]\) is a derivation.

For \eqref{eq:semidirect}, check Jacobi according to the number of
entries in each summand. Three entries from one summand give that
summand's Jacobi identity. Two from \(\mathfrak g\) and one from
\(\mathfrak n\) give
\(\theta([x,y])a-[\theta(x),\theta(y)]a=0\). One from
\(\mathfrak g\) and two from \(\mathfrak n\) give the derivation
identity for \(\theta(x)\). These cases exhaust the expansion.
When \(\mathfrak n\) is abelian, the last condition is automatic,
and the second case is necessary as well as sufficient.
\end{proof}

\section{The universal enveloping algebra}

\begin{theorem}[Construction and universal property]\label{res:A04}
Let \(T(\mathfrak g)=\bigoplus_{r\geq0}\mathfrak g^{\otimes r}\).
The quotient
\[
 U(\mathfrak g)=T(\mathfrak g)/
  \bigl(x\otimes y-y\otimes x-[x,y]\mid x,y\in\mathfrak g\bigr)
\]
has a canonical Lie map \(\iota:\mathfrak g\to U(\mathfrak g)\).
Every Lie map \(\mathfrak g\to A\), with \(A\) associative, extends
uniquely to an algebra map \(U(\mathfrak g)\to A\).
Representations of \(\mathfrak g\) and left \(U(\mathfrak g)\)-modules
are consequently equivalent, with the same linear morphisms.
\end{theorem}

\begin{proof}
The tensor algebra extends a linear map \(f:\mathfrak g\to A\)
uniquely by
\(x_1\otimes\cdots\otimes x_r\mapsto f(x_1)\cdots f(x_r)\).
It annihilates the displayed ideal exactly when
\(f([x,y])=[f(x),f(y)]\). The quotient therefore has the universal
property. Its generators force uniqueness. Taking
\(A=\operatorname{End}_k(M)\) yields the module correspondence.
A linear map commutes with every \(\iota(x)\) if and only if it
commutes with all products and their sums, so the morphisms also
agree. The two constructions leave the underlying action unchanged
and are inverse.
\end{proof}

\begin{theorem}[Poincare--Birkhoff--Witt]\label{res:A05}
Fix a totally ordered basis \((x_i)\) of \(\mathfrak g\). The ordered
monomials
\[
 1,\qquad x_{i_1}\cdots x_{i_r}\quad
 (r\geq1,\ i_1\leq\cdots\leq i_r)
\]
form a basis of \(U(\mathfrak g)\). In particular, \(\iota\) is
injective. For the word-length filtration \(F_rU\),
\[
 S(\mathfrak g)\ \xrightarrow{\ \sim\ }\operatorname{gr}U(\mathfrak g)
\]
is the natural graded algebra isomorphism. In characteristic zero
the symmetrization map
\[
 \operatorname{sym}(y_1\cdots y_r)
 =\frac1{r!}\sum_{\sigma\in S_r}y_{\sigma(1)}\cdots y_{\sigma(r)}
 \tag{1.2}\label{eq:symmetrization}
\]
is a filtered vector-space isomorphism \(S(\mathfrak g)\to U(\mathfrak g)\);
it is equivariant for the adjoint action, but in general is not an
algebra homomorphism.
\end{theorem}

\begin{proof}
In a word, replace a descending adjacent pair \(yx\) by
\(xy+[y,x]\), expressing the bracket in the chosen basis.
Order terms first by word length and then, at fixed length, by
number of inversions. Each replacement produces one word with
one fewer inversion and terms of smaller length. Thus every
reduction terminates.

We verify independence of the first replacement. Disjoint
replacements commute. An overlapping pair can only occur in
a strictly descending triple \(zyx\). Reducing its left pair or
its right pair and then moving the remaining length-three
letters to \(xyz\) leaves a difference of shorter terms
\[
 [y,x]z-z[y,x]+y[z,x]-[z,x]y+[z,y]x-x[z,y].
\]
After reducing each commutator of two basis-linear expressions,
this becomes
\[
 [[y,x],z]+[y,[z,x]]+[[z,y],x]=0.
\]
Induction on length and inversion number therefore gives a unique
normal form. More explicitly, to compare two first reductions,
commute them when disjoint, use this overlap calculation when
adjacent, and normalize all shorter resulting terms by the
induction hypothesis. Applying this comparison successively
compares any two finite reduction sequences.

The normal-form map kills each defining relation with any prefix
and suffix. It consequently kills their linear combinations,
which are exactly the generated two-sided ideal. Conversely,
each reduction changes a word by a defining relation. Thus the
normal forms identify the quotient with the vector space having
the ordered words as basis. In particular, no linear combination
of the length-one words vanishes.

The relation becomes commutativity in top degree, so it defines
the map from \(S(\mathfrak g)\) to \(\operatorname{gr}U\).
Both sides have the same ordered monomial basis in each degree,
which proves bijectivity. In \eqref{eq:symmetrization}, every
permuted product has the same top symbol. Thus symmetrization
induces the identity on the associated graded spaces.
Injectivity and surjectivity now follow by induction on the
finite filtration degree of an element. The formula is
symmetric and multilinear, so is well defined on \(S^r\).
Applying \(\operatorname{ad}x\) to a product differentiates its
factors; averaging commutes with this operation and proves
equivariance.
\end{proof}

This argument works for a free Lie algebra over a commutative
coefficient ring when an ordered module basis is supplied.
The characteristic-zero assumption is needed for
symmetrization, not for the ordered basis theorem. See
\cite[Sections 17.3--17.4]{humphreys} and
\cite[Section 13]{etingof}.

\begin{proposition}[Tensor, dual, and Hopf structures]\label{res:A06}
The enveloping algebra has algebra maps
\[
 \Delta(x)=x\otimes1+1\otimes x,\qquad \epsilon(x)=0
\]
and a linear anti-automorphism \(S\) with \(S(x)=-x\).
They satisfy coassociativity, the counit identities, and
\(m(S\otimes1)\Delta=m(1\otimes S)\Delta=\eta\epsilon\).
For representations \(M,N\),
\[
 x(m\otimes n)=(xm)\otimes n+m\otimes(xn),\qquad
 (x\phi)(m)=-\phi(xm)
 \tag{1.3}\label{eq:tensor-dual-action}
\]
give the tensor and dual representations. On
\(\operatorname{Hom}_k(M,N)\), the corresponding action is
\((x f)(m)=x f(m)-f(xm)\).
\end{proposition}

\begin{proof}
The primitive assignment for \(\Delta\) preserves brackets:
the two tensor factors commute, and the remaining two
commutators are \([x,y]\otimes1+1\otimes[x,y]\).
The universal property constructs \(\Delta\); it also constructs
\(\epsilon\). The assignment \(x\mapsto-x\) is a Lie map to
\(U^{\mathrm{op}}\), so constructs \(S:U\to U^{\mathrm{op}}\).
The two maps \(U\to U^{\otimes3}\) in coassociativity agree on
every generator; the counit identities are checked the same way.
For the antipode identity, the assertion holds on \(1\) and on
each generator. If it holds for \(a,b\), multiplicativity of
\(\Delta\), anti-multiplicativity of \(S\), and the two preceding
identities give it for \(ab\):
\[
 \sum S(a_{(1)}b_{(1)})a_{(2)}b_{(2)}
 =\sum S(b_{(1)})\Big(\sum S(a_{(1)})a_{(2)}\Big)b_{(2)}
 =\epsilon(a)\epsilon(b)1.
\]
Linearity and generation prove it on \(U\).
The first formula in \eqref{eq:tensor-dual-action} is induced
by \(\Delta\), and the second by \(S\).
Alternatively, direct commutator expansion gives the
representation identities and shows that evaluation and
composition are equivariant. Identifying \(M^*\otimes N\)
with \(\operatorname{Hom}_k(M,N)\) when \(M\) is finite
dimensional gives the last formula.
\end{proof}

\section{Free Lie algebras and the BCH product}

\begin{theorem}[The free Lie algebra]\label{res:A07}
For a vector space \(V\) there is a Lie algebra \(L(V)\) and a
linear map \(V\to L(V)\) through which every linear map from \(V\)
to a Lie algebra extends uniquely to a Lie homomorphism.
Moreover,
\[
 U(L(V))\simeq T(V),
\]
and the resulting inclusion identifies \(L(V)\) with the Lie
subalgebra of \(T(V)\) generated by \(V\).
\end{theorem}

\begin{proof}
Take the free nonassociative algebra on \(V\): its degree \(r\)
part is the sum of copies of \(V^{\otimes r}\) indexed by binary
bracketings. Quotient by the ideal generated by alternation
and Jacobi, including all their substitutions and bracket
contexts. Evaluation of a bracketed expression shows the
universal property. No relations on \(V\) are forced, since a
linear map \(V\to V\) with abelian target extends.

For every associative \(A\), a map \(U(L(V))\to A\) is a Lie map
\(L(V)\to A\), hence a linear map \(V\to A\). Construct the
corresponding maps between \(U(L(V))\) and \(T(V)\); both
composites fix their algebra generators \(V\), so are identities.
PBW embeds \(L(V)\) in its enveloping algebra. Its image is
generated under brackets by \(V\), by the nonassociative
construction, proving the last assertion.
\end{proof}

\begin{theorem}[Primitive elements and formal BCH]\label{res:A08}
Suppose \(\operatorname{char}k=0\), and complete \(T(V)\) by tensor
degree. Its elements of positive degree that are primitive for
\(\Delta\) are exactly the degree-completed free Lie algebra
\(\widehat{L(V)}\). Exponentiation and logarithm give mutually
inverse correspondences between primitive elements and
group-like elements with constant term \(1\). Consequently,
\[
 \operatorname{BCH}(X,Y)=\log(\exp X\exp Y)
 =X+Y+\tfrac12[X,Y]
   +\tfrac1{12}[X,[X,Y]]+\tfrac1{12}[Y,[Y,X]]+\cdots
 \tag{1.4}\label{eq:formal-bch}
\]
is a formal Lie series. It satisfies
\(\operatorname{BCH}(\operatorname{BCH}(X,Y),Z)
=\operatorname{BCH}(X,\operatorname{BCH}(Y,Z))\).
\end{theorem}

\begin{proof}
First consider an ordinary, noncompleted enveloping algebra.
Let \(p\) be primitive, and take its highest PBW symbol
\(\bar p\in S^r(L(V))\). The PBW coproduct says that this
symbol is primitive in the symmetric algebra.
In a polynomial algebra over a characteristic-zero field,
the only primitive polynomials are linear: in
\(q(s+t)=q(s)+q(t)\), the coefficient of a single \(t_i\)
is \(\partial q/\partial s_i\), and therefore is constant.
All second derivatives vanish, so \(q\) has degree at most
one. A polynomial uses only finitely many variables, and the
argument applies to it even when the basis is infinite.
Thus \(r\leq1\), so \(p\in k+L(V)\). Comparing constant terms
in the primitive identity removes the scalar.

The tensor grading and the coproduct are compatible. A
primitive element of the completion has each homogeneous
tensor component primitive. Each such component is an
ordinary polynomial in \(T(V)=U(L(V))\), so the preceding
argument puts it in \(L(V)\). Conversely, generators are
primitive, brackets of primitives are primitive, and degreewise
limits preserve the identity. This proves the assertion
about completed primitives.

The series for \(\exp p\) and \(\log(1+a)\) are well defined:
in each tensor degree only finitely many summands contribute.
Their usual inverse identities are formal identities in one
variable. If \(p\) is primitive, then
\[
 \Delta(\exp p)=\exp(p\otimes1+1\otimes p)
              =\exp p\otimes\exp p
\]
because the two summands commute. If \(g\) is group-like,
\(\Delta g=(g\otimes1)(1\otimes g)\), again with commuting
factors, and hence
\(\Delta\log g=\log g\otimes1+1\otimes\log g\).
These computations prove the correspondence.

The product of two group-like elements is group-like.
Therefore its logarithm is a Lie series. Expanding the
associative series through degree three yields the displayed
coefficients in \eqref{eq:formal-bch}. Exponentiating the two
claimed associativity expressions gives
\(\exp X\exp Y\exp Z\) on both sides. Applying the formal
logarithm proves their equality. Substitution is legitimate
because each substituted argument has positive degree.
\end{proof}

\begin{theorem}[Finite BCH on a nilpotent algebra]\label{res:A09}
Let \(\mathfrak n\) be nilpotent of class at most \(c\) over a
characteristic-zero field, so \(\gamma_{c+1}\mathfrak n=0\).
The evaluation of \eqref{eq:formal-bch} is a finite polynomial
product \(x*y\) on \(\mathfrak n\). It is associative, with unit
\(0\) and inverse \(-x\). Lie homomorphisms preserve this
product. If \(D\) is a derivation, then
\[
 D(x*y)=d(\operatorname{BCH})_{(x,y)}(Dx,Dy).
 \tag{1.5}\label{eq:bch-derivation}
\]
\end{theorem}

\begin{proof}
A homogeneous Lie monomial of degree \(r\) is a linear
combination of brackets of length \(r\), which lie in
\(\gamma_r\mathfrak n\) by Proposition~\ref{res:A02}.
Terms of degree greater than \(c\) therefore vanish.
Evaluate the universal formal associativity identity:
all terms of bracket length at most \(c\) form finite sums,
and the others evaluate to zero. This gives associativity.
The formal identities \(\exp0=1\) and
\(\exp x\exp(-x)=1\) give the unit and inverse.

Every Lie homomorphism preserves every finite nested bracket,
so preserves the polynomial. A derivation differentiates a
nested bracket by applying \(D\) to one argument at a time,
by induction on bracket length. Applying this to each term
of the finite BCH polynomial is exactly its polynomial
directional derivative, which proves \eqref{eq:bch-derivation}.
\end{proof}

The construction here is algebraic: it does not use the
integration of an arbitrary Lie algebra into a global Lie
group. That distinction will matter when enlarging a
solvable algebra in the proof of Ado's theorem.
